\section{开题报告}

\subsection{问题提出的背景}

参数化偏微分方程广泛出现于热传导、扩散、渗流、弹性力学、电磁场以及多物理场耦合等问题中。在许多实际应用里，研究者往往并不满足于对某一个固定参数下的模型求解一次，而是需要针对不同参数反复进行数值求解。例如，在参数识别问题中，需要不断调整候选参数并比较模型输出；在优化设计与最优控制中，需要对多个设计变量下的状态方程进行重复评估；在不确定性量化中，需要对参数空间进行大量采样；在实时预测、快速决策和数字孪生等场景中，则要求在参数变化后尽快给出新的近似解。对于这类问题，关键困难不在于单次求解，而在于总体重复求解成本过高，因此常称为 \textbf{many-query 场景}。

以参数化椭圆型方程为例，考虑
\[
-\nabla\cdot(a(x;\mu)\nabla u(x;\mu))=f(x),\qquad x\in\Omega,\ \mu\in\mathcal P\subset\mathbb R^p,
\]
配以适当边界条件。经过有限元离散后，可得到大规模代数系统
\[
A(\mu)u_h(\mu)=f(\mu),\qquad A(\mu)\in\mathbb R^{N\times N},\ u_h(\mu)\in\mathbb R^N,
\]
其中 $N$ 为离散自由度数，通常远大于 $1$。当参数 $\mu$ 在参数域 $\mathcal P$ 中变化时，如果对每一个参数点都独立进行一次高精度有限元求解，那么随着参数样本数的增加，总计算成本将迅速增大。因此，如何在保证近似精度的前提下显著降低重复求解开销，是参数化 PDE 数值计算中的重要问题。

为解决这一问题，模型降阶方法（Model Order Reduction, MOR）得到了广泛研究。其中，\textbf{Reduced Basis（RB）方法} 是参数化 PDE 中最具代表性的降阶框架之一 \cite{hesthaven2015,quarteroni2016,patera_rozza}。其基本思想是：虽然每个高精度解都位于高维空间 $\mathbb R^N$ 中，但对于许多参数化问题，不同参数对应的解往往集中在某个低维流形附近，因此可以构造一个维数远小于 $N$ 的低维子空间来逼近高维解。

若记降阶基矩阵为
\[
V_r\in\mathbb R^{N\times r},\qquad r\ll N,
\]
则可以近似表示
\[
u_h(\mu)\approx V_r a(\mu),
\]
再通过 Galerkin 投影得到低维系统
\[
V_r^T A(\mu)V_r a(\mu)=V_r^T f(\mu).
\]

然而，RB 方法的关键并不只是在于“降维”，更在于如何构造一个尽可能优的低维基空间。最常见的方法是收集若干高精度快照解
\[
U=[u_h(\mu_1),u_h(\mu_2),\dots,u_h(\mu_m)]\in\mathbb R^{N\times m},
\]
并通过奇异值分解
\[
U=\Phi\Sigma\Psi^T
\]
提取前若干个最重要的左奇异向量，作为 POD/RB 基。这种方法在 Frobenius 范数意义下给出最优低秩逼近，因此在模型降阶中具有核心地位 \cite{kunisch2002,trefethen1997}。

进一步的问题是，在 many-query 场景下，快照往往是随着采样过程逐步生成的，而不是一开始就全部获得。如果每增加一个快照都重新对全部快照矩阵做一次 batch SVD，则不仅计算代价高，而且需要较大的存储开销。

基于这一背景，\textbf{增量奇异值分解（incremental SVD）} 成为一个自然且重要的研究方向。Brand 在文献 \cite{brand2002,brand2006} 中提出了经典的增量 SVD 更新方法，可以在新数据到达时更新奇异值分解，而无需重新分解整个矩阵。

\subsubsection{背景介绍}

从数值分析角度看，参数化 PDE 的 many-query 计算具有典型的“离线贵、在线快”的需求结构。RB 方法通常采用 Offline/Online 分解：离线阶段生成快照、构造基空间；在线阶段仅在低维空间内求解，从而大幅减少计算量 \cite{hesthaven2015}。

设快照矩阵为
\[
U=[u_1,u_2,\dots,u_m],
\]
经典方法需要对其统一执行 SVD 分解。而在实际构基中，快照往往逐步生成，因此矩阵具有按列扩展的结构。

在 weighted inner product 下，若
\[
U_\ell = Q\Sigma R^T,\qquad Q^TWQ=I,
\]
则称其为 core SVD。Brand 型增量算法通过将新增快照分解为投影部分和残差部分：

\[
d=Q^T(Wu_{\ell+1}),\qquad
e=u_{\ell+1}-Qd,
\qquad
p=(e^TWe)^{1/2},
\]

从而构造更新矩阵

\[
Y=
\begin{bmatrix}
\Sigma & d\\
0 & p
\end{bmatrix}.
\]

对该小矩阵执行 SVD，即可完成整体更新。

从数学角度看，POD 方法与奇异值分解之间存在紧密联系。
设快照矩阵为
\[
U=[u_1,u_2,\dots,u_m]\in\mathbb R^{N\times m},
\]
其奇异值分解为
\[
U=\Phi\Sigma\Psi^T,
\]
其中 $\Phi$ 和 $\Psi$ 分别为左右奇异向量矩阵，
$\Sigma$ 为奇异值矩阵。
经典结果表明，
取前 $r$ 个左奇异向量所张成的子空间
在 Frobenius 范数意义下
给出了对快照矩阵的最优秩-$r$ 逼近
\cite{trefethen1997}。
因此，
POD 基可以被理解为
快照数据能量最大的主导模态。

然而，
在参数化 PDE 的实际计算中，
快照数量往往不断增加，
如果每一次新增快照都重新执行完整 SVD，
则离线阶段计算成本将迅速增长。
因此，
如何在已有 SVD 基础上进行高效更新，
成为模型降阶中的重要问题。
增量奇异值分解正是在这一背景下提出，
其目标是在新数据到达时更新奇异值分解，
而不必重新处理全部数据。


近年来，增量 POD 方法也被用于 PDE 仿真数据处理中 \cite{fareed2018,fareed2020}。此外，Zhang 在 \cite{zhang2022} 中研究了 weighted inner product 下的增量 SVD 算法，并分析了其数值稳定性。

\subsubsection{POD 方法的局限性与增量 POD 的必要性}
POD 方法由于能够从快照数据中提取主导模态，并在 Frobenius 范数意义下给出最优低秩逼近，因此已成为模型降阶中最常用的构基方法之一。然而，经典 POD 方法通常建立在 \textbf{batch SVD} 的基础上，即必须在收集全部快照之后，对完整快照矩阵一次性执行奇异值分解。这样的处理方式在理论上是清晰的，但在参数化 PDE 的实际计算中存在若干局限性。

首先，\textbf{经典 POD 的计算代价较高}。设快照矩阵为
\[
U=[u_1,u_2,\dots,u_m]\in\mathbb R^{N\times m},
\]
其中 $N$ 是有限元离散后的空间维数，$m$ 是快照数量。在参数化 PDE 的 many-query 场景下，$N$ 往往很大，而 $m$ 也会随着参数采样不断增加。此时，若始终采用 batch SVD，则需要反复对大规模矩阵进行分解，离线阶段的计算成本会迅速升高。尤其当每增加一个新快照都重新对整个快照矩阵做一次 SVD 时，这种重复计算会显著削弱模型降阶在整体流程中的效率优势。

其次，\textbf{经典 POD 对快照存储的要求较高}。batch SVD 通常要求在离线阶段保存全部历史快照，以便统一构造快照矩阵并进行分解。然而，在有限元背景下，每个快照本身就是一个高维向量，当快照数量逐步增加时，存储完整矩阵会带来明显的内存压力。对于大规模参数扫描、长时间仿真或流式生成数据的场景，这一问题尤为突出。

再次，\textbf{经典 POD 不适合快照逐步到达的构基过程}。在实际 reduced basis 方法中，快照往往不是预先全部给定的，而是随着贪婪选点、自适应采样或参数扫描过程逐步生成的。换言之，快照矩阵在离线阶段常常具有“按列不断扩展”的结构。如果仍坚持每次新增样本后都重新执行完整 SVD，则方法在算法结构上与实际数据生成机制并不匹配。因此，经典 POD 更适合静态、一次性的数据处理，而不够适合逐步更新的离线构基过程。

基于以上原因，研究\textbf{增量 POD}方法具有明显必要性。所谓增量 POD，是指在新快照到达时，不必重新处理全部历史数据，而是在已有分解结果的基础上递推更新当前低维基空间。本文中所讨论的增量 POD，主要通过对快照矩阵进行\textbf{增量奇异值分解（incremental SVD）}来实现。其核心思想是：当新的快照列向量加入时，将其分解为相对于当前基空间的投影部分与正交残差部分，再通过一个小规模矩阵的奇异值分解完成整体更新。这样就避免了对整个扩展后快照矩阵重复做 batch SVD。

与经典 POD 相比，增量 POD 主要具有以下几个方面的优势。第一，它能够\textbf{适应快照逐步生成的离线流程}，与参数化 PDE 中常见的贪婪采样、自适应选点等策略更加匹配。第二，它能够\textbf{减少重复分解带来的计算负担}，从而提高离线构基效率。第三，它能够\textbf{降低对完整历史快照存储的依赖}，因而在大规模问题中更具可扩展性。第四，它使得“每增加一个样本就同步更新降阶基”成为可能，从而为更灵活的构基方式提供了算法基础。

当然，增量 POD 并不是对经典 POD 的简单替代。与 batch POD 相比，它虽然在效率和存储方面更具优势，但也带来了新的数值问题。最典型的问题是：在多次递推更新过程中，浮点误差会逐渐积累，导致基向量的正交性退化，进而影响奇异值的准确性、截断策略的可靠性以及最终降阶基的质量。特别是在有限元离散诱导的 weighted inner product 框架下，这一问题会更加明显。因此，增量 POD 的研究重点不仅在于“如何更快地更新”，也在于“如何在更新过程中保持数值稳定性”。

\subsection{论文的主要内容和技术路线}

\subsubsection{主要研究内容}

本文拟以参数化椭圆型偏微分方程为研究对象，围绕“增量奇异值分解在 PDE 模型降阶中的应用”展开研究，主要内容包括以下几个方面。

第一，建立参数化 PDE 模型降阶的基本框架。本文将从参数化椭圆型方程出发，给出其有限元离散形式，并说明在 many-query 场景下直接求解高维系统的成本问题。在此基础上，引入 Reduced Basis 方法和 POD 构基思想，说明快照矩阵的低秩结构为何能够支持模型降阶，以及 SVD 在构造近似最优低维基中的作用。

第二，分析 weighted inner product 下的 core SVD 结构。由于有限元离散通常伴随质量矩阵，因此本文将重点讨论加权内积
\[
\langle u,v\rangle_W=u^TWv
\]
下的正交性与矩阵分解形式，解释为何在 PDE 快照处理中应当考虑 core SVD，而不是简单套用标准 Euclidean SVD。这里的目标不是建立全新的抽象理论，而是在已有定义基础上，准确梳理其与有限元离散和 POD 构基之间的关系。

第三，研究增量 SVD 的初始化与更新机制。设快照矩阵按列逐步生成，本文将分析在已知
\[
U_\ell = Q\Sigma R^T
\]
的条件下，如何通过新增列向量 $u_{\ell+1}$ 构造
\[
U_{\ell+1}=[U_\ell,u_{\ell+1}]
\]
的更新分解。重点包括：投影项与残差项的分离、小规模边界矩阵 $Y$ 的构造、更新后左右奇异向量与奇异值的表达方式，以及当残差很小或新奇异值过小时的截断处理。

第四，讨论重正交与数值稳定性问题。增量 SVD 的主要数值风险在于正交性退化，因此本文将分析为什么浮点误差会导致
\[
Q^TWQ\neq I,
\]
以及这种退化对奇异值真实性、截断可靠性和降阶基质量的影响。在此基础上，讨论加权 Gram--Schmidt 型重正交化的作用及其计算代价，并尝试比较不同策略下的数值表现。

第五，设计数值实验并进行结果比较。本文计划选取一个或若干典型的参数化椭圆型方程作为实验对象，生成有限元快照数据，并分别采用 batch SVD 与增量 core SVD 构造 POD/RB 基。比较指标包括：奇异值衰减、降阶逼近误差、加权正交性误差、离线时间开销以及内存使用情况。通过这些结果，评估增量 SVD 在 PDE 模型降阶中的适用范围与优势。

\subsubsection{技术路线}

本文拟采用“理论梳理—算法实现—数值比较”的技术路线。

第一步，完成文献阅读与问题建模。首先系统阅读参数化 PDE 模型降阶、RB/POD 方法、增量 SVD 以及 weighted inner product 下 core SVD 相关文献，明确课题中的数学对象、算法目标与关键困难。随后选定一个典型的参数化椭圆型方程作为主要测试模型，并确定其边界条件、参数形式和有限元离散方案。

第二步，建立高精度快照生成模块。对选定的参数化 PDE 进行有限元离散，得到高维代数系统
\[
A(\mu)u_h(\mu)=f(\mu).
\]
在给定参数采样集 $\{\mu_1,\dots,\mu_m\}$ 下，计算对应的高精度数值解，并将其组织为快照矩阵
\[
U=[u_h(\mu_1),\dots,u_h(\mu_m)].
\]
同时构造相应的权矩阵 $W$，通常取有限元质量矩阵，以便后续进行 weighted core SVD。

第三步，实现 batch SVD 与增量 SVD 两套离线构基流程。对于 batch 方法，直接对快照矩阵做标准或加权意义下的分解并截断，得到 POD/RB 基；对于增量方法，则按照“初始化—更新—必要时重正交—截断”的流程，逐列处理快照矩阵。这里的重点是保证实现与数学公式一致，并对关键中间量如残差范数 $p$、奇异值阈值和正交性误差进行记录。

第四步，构造降阶模型并进行在线测试。将所得到的低维基代入参数化椭圆方程的离散系统，建立降阶模型
\[
V_r^T A(\mu)V_r a(\mu)=V_r^T f(\mu),
\]
对若干测试参数进行在线求解，并将降阶解与高精度有限元解进行比较，从而评估不同构基方式对在线精度的影响。

第五步，分析实验结果并总结结论。比较 batch SVD 与增量 SVD 在离线时间、内存开销、正交性保持、降阶误差和基维数控制等方面的差异，分析不同截断阈值与重正交策略的影响，并据此总结增量 SVD 在参数化 PDE 模型降阶中的优缺点与适用条件。

\subsubsection{可行性分析}

从研究对象看，本文拟研究的参数化椭圆型方程属于数值分析与科学计算中较为成熟的一类模型，其有限元离散、矩阵组装与数值求解流程较为清晰，适合作为本科阶段模型降阶研究的基础载体。因此，在问题设置上具有较强可行性。

从理论基础看，课题涉及的核心概念包括有限元方法、POD/RB 模型降阶、奇异值分解与 weighted inner product 下的 core SVD。这些内容虽然跨越 PDE 数值、矩阵计算和模型降阶三个方向，但彼此联系紧密，且已有较为清楚的文献基础。本文的目标不是从零建立全新的深层理论，而是在已有理论框架下完成准确梳理、合理实现与数值分析，因此对于本科毕业论文而言是合适的。

从技术实现看，本课题所需的主要计算步骤包括：参数化 PDE 的有限元求解、快照矩阵构造、SVD 或增量 SVD 实现、重正交处理以及降阶误差比较。这些步骤均可借助现有数值计算环境完成，例如 MATLAB、Python 或其他常用科学计算工具。特别是增量 SVD 的更新公式本身由小矩阵 SVD 和向量投影构成，实现难度可控，便于逐步调试和验证。

从工作量看，若将研究重点聚焦在“参数化椭圆方程 + weighted incremental core SVD + 数值比较”这一主线上，则论文内容既有足够的数学深度，又不至于因为问题过大而难以完成。相比试图证明过于复杂的理论结果，这种“数学框架清晰、算法流程完整、实验结果可展示”的选题更符合本科毕业论文的时间与能力边界。

此外，从算法复杂度角度看，
增量 SVD 方法在处理大规模快照矩阵时
具有明显优势。
传统 batch SVD 的计算复杂度通常为
$\mathcal O(Nm^2)$ 或 $\mathcal O(N^2m)$，
当快照维数 $N$ 和样本数量 $m$ 都较大时，
计算成本十分昂贵。
相比之下，
增量 SVD 每次只需要对一个
$(r+1)\times(r+1)$ 的小矩阵执行 SVD，
其计算复杂度显著降低。
因此，
在 many-query 场景下，
采用增量 SVD 构造降阶基
不仅可以减少计算时间，
还可以降低内存存储需求，
从而提高离线阶段的整体效率。

\subsection{研究计划进度安排及预期目标}

\subsubsection{进度安排}

第一阶段：文献调研与问题细化。系统阅读参数化 PDE、Reduced Basis 方法、POD/SVD、incremental SVD 以及 weighted core SVD 相关文献，梳理课题中的基本概念、主要算法与已有结果。在此基础上，明确本文所采用的参数化椭圆方程模型、有限元离散方式以及数值实验平台。

第二阶段：建立高精度模型与快照生成程序。完成参数化椭圆型方程的有限元离散与求解程序，测试不同参数下高精度数值解的正确性与稳定性。随后生成训练快照集与测试快照集，为后续 batch SVD 和增量 SVD 构基提供数据基础。

第三阶段：实现 batch SVD 与增量 core SVD 算法。首先实现标准的快照矩阵分解与 POD 基提取过程，作为比较基线；再实现 weighted inner product 下的增量更新、截断与重正交步骤，并逐步验证更新公式是否正确、正交性是否保持以及截断策略是否有效。

第四阶段：建立降阶模型并开展数值实验。基于两种构基方式得到的低维基，建立参数化椭圆方程的降阶模型，比较其在不同参数点上的近似误差、离线耗时、内存消耗与正交性误差。对实验结果进行整理，并尝试分析不同阈值、不同基维数下的表现差异。

第五阶段：整理结论并撰写论文。对已有理论梳理、算法流程和实验结果进行系统总结，完成文献综述、开题报告、正文初稿与修改稿的撰写，并根据指导教师意见进一步修订，最终形成完整毕业论文。

\subsubsection{预期目标}

第一，较系统地梳理参数化椭圆型 PDE 在 many-query 场景下面临的计算问题，说明 Reduced Basis 方法、POD/SVD 以及增量 SVD 之间的逻辑关系，形成较为清晰的理论背景与问题表述。

第二，在 weighted inner product 框架下，较准确地给出增量 core SVD 的数学描述，包括初始化、更新、截断和重正交等步骤，并解释其在有限元快照数据中的意义。

第三，完成一个可运行的数值实验流程：从高精度有限元快照生成，到 batch SVD 与增量 SVD 构基，再到降阶模型测试与结果比较。通过实验展示增量 SVD 在离线构基中的潜在优势，以及其数值稳定性方面需要关注的问题。

第四，形成一篇结构完整、数学表述准确、实验结果明确的本科毕业论文。论文不追求过度扩张结论，而是力求在选定模型和算法框架下，给出可信、清楚、可复现的分析与比较结果，为后续进一步研究增量 SVD 在更复杂 PDE 模型中的应用打下基础。


% 参考文献
{
\newpage
\bibliographystyle{gbt7714-numerical}
\phantomsection
\subsection{参考文献}
{\normalfont\CJKfamily{simsun}\zihao{5}\setlength{\baselineskip}{14pt}
\renewcommand{\refname}{\vspace{-\baselineskip}}
\bibliography{refs/02_Report}}
}